%% P4 — Singapore
%% Target journal: Management International Review (MIR), Springer
%% Document class: svjour3 (Springer journal; svjour3.cls required)
%% Fallback: standard article class if svjour3.cls unavailable
%% Compile: pdflatex → biber → pdflatex × 2  |  or: latexmk -pdf

%% NOTE: For final submission, download svjour3.cls from
%%   https://www.springernature.com/gp/authors/campaigns/latex-author-support
%%   and place in same directory as this .tex file.

\documentclass[smallextended]{svjour3}   % Springer MIR class
% \documentclass[12pt,a4paper]{article}   % fallback if svjour3 unavailable

\smartqed  % Springer: flush right QED boxes

%% ---- Packages ----
\usepackage{amsmath,amssymb,amsthm}
\usepackage{graphicx}
\usepackage{booktabs}
\usepackage[style=apa,backend=biber]{biblatex}
\addbibresource{../../references/p4_references.bib}
\usepackage{hyperref}
\hypersetup{colorlinks=true,linkcolor=black,citecolor=black,urlcolor=blue}

%% ---- Shared macros ----
%% Shared macros for PhD dissertation papers
%% Internationalization-Performance relationship studies (P3, P4, P5, P6)

%% --- Key variables (italic, math mode) ---
\newcommand{\lnLP}{\ln(\mathrm{LP})}
\newcommand{\FSTS}{FSTS}
\newcommand{\FSTSc}{FSTS_{c}}
\newcommand{\FSTScsq}{FSTS_{c}^{2}}
\newcommand{\FSTSccb}{FSTS_{c}^{3}}
\newcommand{\TCI}{\mathrm{TCI}}
\newcommand{\TCIz}{\mathrm{TCI}_{z}}
\newcommand{\DAI}{\mathrm{DAI}}
\newcommand{\DAIz}{\mathrm{DAI}_{z}}
\newcommand{\IMR}{\mathrm{IMR}}

%% --- Model notation ---
\newcommand{\bvec}[1]{\boldsymbol{#1}}
\newcommand{\bX}{\bvec{X}}
\newcommand{\bgamma}{\bvec{\gamma}}
\newcommand{\bcontrols}{\bgamma'\bX_{i}}

%% --- ICRV regimes ---
\newcommand{\ICRVadv}{\mathrm{ICRV}_{\mathrm{Adv}}}
\newcommand{\ICRVem}{\mathrm{ICRV}_{\mathrm{Em}}}
\newcommand{\ICRVfr}{\mathrm{ICRV}_{\mathrm{Fr}}}

%% --- Fixed effects ---
\newcommand{\alphaj}{\alpha_{j}}   % country FE
\newcommand{\deltat}{\delta_{t}}   % year FE

%% --- Turning point shorthand ---
\newcommand{\TP}{\mathrm{TP}}

%% --- Statistical operators ---
\newcommand{\pval}{p}
\newcommand{\betacoef}{\hat{\beta}}


%% ---- Journal info ----
\journalname{Management International Review}

% ============================================================
\begin{document}

\title{Technological Capability, Digital Adoption, and the
  Internationalization--Performance Relationship: A Firm-Level Study of
  Singapore}

\author{[Blinded for review]}

\maketitle

%% ---- ABSTRACT ----
\begin{abstract}
This study examines how technological capability (\(\TCI\)) and digital
adoption (\(\DAI\)) are associated with the internationalization--performance
(\(I \to P\)) relationship among firms in Singapore --- a boundary-case test
at the advanced-economy ceiling of the ICRV spectrum, where digital
infrastructure is near-universally diffused.

Using World Bank Enterprise Survey microdata (Singapore 2023, \(N = 623\)),
the analysis separates firm-internal capability depth (\(\TCI\)) from
foundational digital interfaces and transaction-enabling mechanisms (\(\DAI\)).

Three findings emerge.  First, the \(I \to P\) relationship within the
observed export-intensity range is better characterised as predominantly
positive with mild quadratic curvature than as a classic inverted-\(U\); the
implied turning point lies in a sparsely populated upper tail (support-constrained
estimate: 53--253\% FSTS, Lind--Mehlum \(p = .303\)).  Second, \(\TCI\) is
positively associated with labour productivity, but no \(\TCI\) moderation of
the \(I \to P\) curve is detected.  Third, \(\DAI\) does not produce a uniform
productivity premium; instead, the \(\FSTScsq \times \DAI\) interaction
coefficient is positive and significant (\(\hat{\beta} = +3.119\),
\(p = .012\)), consistent with digital adoption acting as a contingent scaling
resource that activates under high cross-border coordination demands.

\keywords{internationalization--performance \and digital adoption \and
  technological capability \and export intensity \and labour productivity \and
  Singapore}
\end{abstract}

% ============================================================
\section{Introduction}
\label{sec:intro}

Singapore occupies an analytically distinctive position in the
internationalization--performance (\(I \to P\)) literature.  As an
Advanced~I economy in the ICRV framework, Singapore hosts firms operating at
the digital-infrastructure ceiling --- near-universal internet penetration,
extensive e-payment systems, and cross-border digital trade facilitation.  This
context allows a boundary-case test: does foundational digital adoption still
moderate the \(I \to P\) relationship when it is near-universally available?

% ============================================================
\section{Theory and Hypotheses}
\label{sec:theory}

\subsection{Inverted-\textit{U} and Positive-Saturation (H1)}

\begin{quote}
\textbf{H1:} In the Singapore context (ICRV Advanced~I), the positive
\(I \to P\) relationship may reach a positive-saturation zone rather than a
classic interior maximum, given the support-constrained export-intensity
distribution.
\end{quote}

\subsection{TCI Direct Effect Without Moderation (H2)}

\begin{quote}
\textbf{H2:} \(\TCI\) raises the productivity level (direct positive effect),
but does not significantly moderate the \(I \to P\) curvature in an
advanced-economy context where institutional infrastructure is uniformly high.
\end{quote}

\subsection{DAI as Contingent Scaling Resource (H3)}

\begin{quote}
\textbf{H3:} \(\DAI\) moderates the quadratic \(I \to P\) curve positively at
high export intensity: the \(\FSTScsq \times \DAI\) interaction is positive,
reflecting digital adoption as a coordination-demand amplifier.
\end{quote}

% ============================================================
\section{Methodology}
\label{sec:methods}

\subsection{Data}

WBES Singapore 2023, \(N = 623\) (all firms); exporters sub-sample
\(n = 101\).  All monetary values in 2017 constant USD PPP.

\subsection{Variables}

\paragraph{Dependent.}
\(\lnLP = \ln(\text{annual revenue PPP} / \text{permanent workers})\).

\paragraph{Internationalisation.}
\(\FSTSc = \FSTS - \bar{F}\) (mean-centred to reduce collinearity with
\(\FSTScsq\)).

\paragraph{Technological Capability Index.}
\begin{equation}
  \TCIz = \frac{1}{3}\bigl[z(\text{quality\_cert}) + z(\text{for\_lic\_tech})
    + z(\text{R\&D\_spend})\bigr].
  \label{eq:tci_p4}
\end{equation}

\paragraph{Digital Adoption Index.}
Tier-1 + Tier-2 composite:
\begin{equation}
  \DAI = \frac{1}{3}\bigl[\text{website} + \text{e\_payment\_k33}
    + \text{e\_payment\_k38}\bigr].
  \label{eq:dai_p4}
\end{equation}

\subsection{Model Specification}

\paragraph{M0 --- Linear baseline.}
\begin{equation}
  \lnLP_{i} = \beta_{0} + \beta_{1}\,\FSTS_{i}
    + \bcontrols + \alphaj + \varepsilon_{i}
  \label{eq:M0_p4}
\end{equation}

\paragraph{M1 --- Quadratic (H1).}
\begin{equation}
  \lnLP_{i} = \beta_{0} + \beta_{1}\,\FSTSc_{i} + \beta_{2}\,\FSTScsq_{i}
    + \bcontrols + \alphaj + \varepsilon_{i}
  \label{eq:M1_p4}
\end{equation}
Turning point: \(\TP = -\beta_{1}/(2\beta_{2})\).

\paragraph{M2 --- TCI moderation (H2).}
\begin{align}
  \lnLP_{i} &= \beta_{0} + \beta_{1}\,\FSTSc_{i} + \beta_{2}\,\FSTScsq_{i}
    \notag \\
    &\quad + \beta_{3}\,\TCIz_{i}
    + \beta_{4}\,(\FSTSc_{i} \times \TCIz_{i})
    + \beta_{5}\,(\FSTScsq_{i} \times \TCIz_{i})
    \notag \\
    &\quad + \bcontrols + \alphaj + \varepsilon_{i}
  \label{eq:M2_p4}
\end{align}

\paragraph{M3 --- DAI moderation (H3).}
\begin{align}
  \lnLP_{i} &= \beta_{0} + \beta_{1}\,\FSTSc_{i} + \beta_{2}\,\FSTScsq_{i}
    \notag \\
    &\quad + \beta_{3}\,\DAI_{i}
    + \beta_{4}\,(\FSTSc_{i} \times \DAI_{i})
    + \beta_{5}\,(\FSTScsq_{i} \times \DAI_{i})
    \notag \\
    &\quad + \bcontrols + \alphaj + \varepsilon_{i}
  \label{eq:M3_p4}
\end{align}
Key test: \(\hat{\beta}_{5} > 0\) (positive quadratic amplification).
Observed: \(\hat{\beta}_{5} = +3.119\), \(p = .012\).

\paragraph{M4 --- Heckman two-step IMR check.}
To address exporter selection bias, a Heckman first-stage probit models
\(\Pr(\FSTS > 0)\) on firm size, age, foreign ownership, and industry.
The Inverse Mills Ratio (\(\IMR\)) is added as a regressor in M4:
\begin{equation}
  \lnLP_{i} = \beta_{0} + \beta_{1}\,\FSTSc_{i} + \beta_{2}\,\FSTScsq_{i}
    + \beta_{3}\,\DAI_{i} + \beta_{4}\,(\FSTScsq_{i} \times \DAI_{i})
    + \lambda\,\IMR_{i} + \bcontrols + \alphaj + \varepsilon_{i}
  \label{eq:M4_p4}
\end{equation}

\paragraph{Controls.}
\(\bX_{i}\): \(\ln(\text{employees})\), firm age, foreign-ownership dummy,
sector dummies; \(\alphaj\) = industry FE (country FE collinear for single
country).  HC1 robust standard errors.

\paragraph{Statistical note.}
Exporters sub-sample (\(n = 101\)) yields power \(\approx 16\%\) for
\(f^{2} = 0.02\) moderator --- an inferential bound that should be interpreted
cautiously.  The positive-saturation turning-point estimate (53--253\%) is
support-constrained and reported as exploratory.

% ============================================================
\section{Results}
\label{sec:results}

\subsection{Main Effects}

\(\TCIz\) raises \(\lnLP\) in the full sample (\(\hat{\beta} = 0.283\),
\(p < .001\)).  No significant \(\TCI\) moderation of the \(I \to P\) curve
is detected (joint \(F\)-test \(p = .412\)), consistent with H2.

\subsection{DAI Conditional Amplification}

The \(\FSTScsq \times \DAI\) coefficient is \(+3.119\) (\(p = .012\)),
indicating that digital adoption amplifies productivity at high export intensity.
The Lind--Mehlum test (p = .303) signals positive saturation rather than a
classic interior maximum within the observed FSTS range.

\subsection{Heckman Correction}

The IMR coefficient is not statistically significant (\(p = .318\)),
suggesting exporter selection bias does not materially confound the main
estimates.

% ============================================================
\section{Discussion and Conclusion}
\label{sec:discussion}

The findings from Singapore confirm the \emph{Context-Contingent Digital
Capability Model} (CDCM): in an Advanced~I ICRV economy, foundational digital
adoption functions as a contingent scaling resource rather than a universal
productivity driver.  The null \(\TCI\) moderation result is consistent with
the ICRV institutional homogeneity argument --- when all firms operate in a
high-quality institutional environment, technological capability raises
productivity levels but does not differentially amplify internationalisation
returns.

\paragraph{Limitations.}
Power constraints in the exporters sub-sample (\(n = 101\)) limit inferential
precision for moderation tests.  The turning-point confidence interval
(53--253\%) is a support-constrained estimate; readers should treat it as a
positive-saturation signal rather than a precisely estimated maximum.

\printbibliography

\end{document}
